\section{实验评估}
\label{sec:evaluation}

本节评估 QMAP 在真实程序访存 trace 上的页面迁移效果。与早期基于合成 trace 的原型验证不同，本文实验主线使用 PARSEC benchmark 采集得到的真实访存序列，并以标准时间顺序切分、压力窗口、真实消融和多随机种子稳定性共同构成证据链。为保持论文表述一致，后文统一将最终方法称为 QMAP。QMAP 的模型结构为 Transformer Encoder 对访问历史编码，然后在时间维上进行均值池化，并结合候选页特征输出迁移优先级。

实验主要回答四个问题。第一，QMAP 在真实 1M trace 上是否能够降低混合内存系统的加权访问代价。第二，当标准时间切分缺少足够 eviction 压力时，QMAP 在压力窗口中是否仍然有效。第三，读写类型特征和代价感知训练目标是否对真实 workload 有贡献。第四，QMAP 的正结果和负结果是否对训练随机种子稳定。

\subsection{实验设置}

\textbf{系统模型。}
我们使用 trace-driven replay 模拟 DRAM/NVM 混合内存系统。DRAM 是容量受限的快速层，NVM 是容量更大的慢速层。当一次访问未命中 DRAM 且 DRAM 已满时，替换策略需要从候选页面集合中选择一个页面迁出 DRAM。QMAP 将这一过程建模为候选页面排序问题：模型读取最近的访存历史、读写类型、程序上下文和候选页统计特征，为每个候选页面输出迁移分数，并选择分数最高的页面迁移到 NVM。

\textbf{数据集。}
真实 workload 来自 PARSEC benchmark，包括 \texttt{blackscholes}、\texttt{canneal}、\texttt{streamcluster} 和 \texttt{dedup}。我们使用 DynamoRIO/drmemtrace 采集用户态访存记录，并转换为 QMAP 使用的 \texttt{PC,Address,RW} 格式。每个 workload 采集 1M 条访存记录，并按照时间顺序切分为 80\% 训练、10\% 验证和 10\% 测试。表~\ref{tab:real-trace-stats} 给出 trace 统计。可以看到，不同 workload 在 unique pages、unique PCs 和写比例上差异明显：\texttt{blackscholes} 的工作集很小，\texttt{streamcluster} 和 \texttt{dedup} 的页面数更多，\texttt{dedup} 的写比例最高。

\begin{table}[t]
  \centering
  \caption{1M PARSEC 真实 trace 统计。所有 workload 均采用 800K/100K/100K 的时间顺序切分。}
  \label{tab:real-trace-stats}
  \scriptsize
  \begin{tabular}{lrrrr}
    \toprule
    Workload & Records & Unique pages & Unique PCs & Write ratio \\
    \midrule
    blackscholes & 1,000,000 & 23 & 611 & 0.3468 \\
    canneal & 1,000,000 & 254 & 6,821 & 0.2854 \\
    streamcluster & 1,000,000 & 767 & 4,238 & 0.3495 \\
    dedup & 1,000,000 & 493 & 34 & 0.5000 \\
    \bottomrule
  \end{tabular}
\end{table}

\textbf{比较策略和参数。}
我们将 QMAP 与 LRU、Random、LFU 和 CLOCK 比较。LRU 和 CLOCK 代表常见低开销替换策略，LFU 代表频率统计基线，Random 用于给出无状态下界。除特别说明外，QMAP 使用相同的默认配置：历史长度为 10，候选页数量为 8，DRAM 容量为 16 页，lookahead 为 256，训练 10 个 epoch，batch size 为 32，随机种子为 3136859。

\textbf{指标。}
我们报告 DRAM 命中率、NVM 写入次数、加权访问代价、迁移次数和平均决策延迟。加权访问代价是主要指标，其代价模型为：DRAM read/write 均为 1，NVM read 为 2，NVM write 为 8，页面迁移为 10。该指标比命中率更适合 DRAM/NVM 混合内存，因为一次 NVM 写入或页面迁移的代价显著高于 DRAM 命中。决策延迟来自 Python replay 原型，主要用于比较相对在线开销，并不代表优化后的系统实现延迟。

\subsection{1M 标准切分主实验}

表~\ref{tab:real-main} 展示 1M 标准时间切分下的端到端结果。QMAP 在 \texttt{blackscholes} 上相较最佳 baseline 将加权访问代价降低 0.91\%，在 \texttt{dedup} 上与 LRU 持平，在 \texttt{streamcluster} 标准测试段上所有策略都没有触发迁移，因而该标准切分不能形成有效策略比较。与此同时，QMAP 在 \texttt{canneal} 上明显失败，相较最佳 baseline 代价升高 19.32\%。这一组结果说明，真实 trace 中 workload 压力差异很大，不能仅凭标准切分得到单一结论：QMAP 在小工作集或低压力场景中可能与传统策略接近，在候选页排序出现偏差时也会产生明显负例。

\begin{table*}[t]
  \centering
  \caption{1M 标准时间切分下的主实验结果。Cost delta 表示 QMAP 相对该 workload 最佳 baseline 的加权代价变化，负数表示 QMAP 更好。}
  \label{tab:real-main}
  \scriptsize
  \begin{tabular}{llrrrrr}
    \toprule
    Workload & Policy & Hit rate(\%) & NVM writes & Cost & Migrations & Decision ms \\
    \midrule
    blackscholes & LRU & 98.89 & 144 & 112,958 & 1,098 & 0.000288 \\
    blackscholes & Random & 98.66 & 145 & 115,505 & 1,329 & 0.001238 \\
    blackscholes & LFU & 99.47 & 221 & 106,952 & 510 & 0.005870 \\
    blackscholes & CLOCK & 99.09 & 107 & 110,437 & 889 & 0.002007 \\
    blackscholes & QMAP & 99.46 & \textbf{32} & \textbf{105,983} & 525 & 3.387825 \\
    \midrule
    canneal & LRU & \textbf{97.63} & 52 & \textbf{126,178} & \textbf{2,350} & 0.000287 \\
    canneal & Random & 96.50 & 193 & 139,465 & 3,481 & 0.001163 \\
    canneal & LFU & 94.69 & 280 & 159,919 & 5,293 & 0.005749 \\
    canneal & CLOCK & 97.62 & 60 & 126,325 & 2,359 & 0.001727 \\
    canneal & QMAP & 95.42 & \textbf{51} & 150,559 & 4,567 & 2.401794 \\
    \midrule
    streamcluster & LRU & 99.99 & 4 & 100,033 & 0 & 0.000000 \\
    streamcluster & Random & 99.99 & 4 & 100,033 & 0 & 0.000000 \\
    streamcluster & LFU & 99.99 & 4 & 100,033 & 0 & 0.000000 \\
    streamcluster & CLOCK & 99.99 & 4 & 100,033 & 0 & 0.000000 \\
    streamcluster & QMAP & 99.99 & 4 & 100,033 & 0 & 0.000000 \\
    \midrule
    dedup & LRU & \textbf{99.95} & \textbf{50} & \textbf{100,734} & \textbf{38} & 0.000572 \\
    dedup & Random & 99.93 & 56 & 100,968 & 56 & 0.001950 \\
    dedup & LFU & 99.89 & 106 & 101,686 & 94 & 0.006607 \\
    dedup & CLOCK & 99.94 & 51 & 100,751 & 39 & 0.002321 \\
    dedup & QMAP & \textbf{99.95} & \textbf{50} & \textbf{100,734} & \textbf{38} & 18.408344 \\
    \bottomrule
  \end{tabular}
\end{table*}

\texttt{blackscholes} 是标准切分中的稳定正结果。虽然它的 active page set 较小，导致所有策略命中率都很高，但 QMAP 仍将 NVM 写入从 LFU 的 221 次降至 32 次，并将加权访问代价从最佳 baseline 的 106,952 降至 105,983。这说明 QMAP 的收益不一定来自显著提升命中率，而可能来自更好地规避高代价写入和迁移。

\texttt{canneal} 是明确的负例。QMAP 的 NVM 写入数与 LRU 接近，分别为 51 和 52，但迁移次数从 LRU 的 2,350 增加到 4,567，最终导致加权访问代价显著升高。该结果表明，QMAP 在该 workload 上不是写入判断失败，而是候选迁移顺序不够稳健，导致过度迁移。这一负例对论文同样重要，因为它界定了 QMAP 当前版本的适用边界。

\texttt{streamcluster} 标准测试段没有触发 eviction 决策，所有策略 cost 完全相同。因此该标准切分不能用于评价替换策略优劣。为避免用低压力测试段掩盖策略差异，我们进一步构造 pressure window，专门评估存在大量迁移决策的阶段。

\subsection{压力窗口实验}

表~\ref{tab:pressure-window} 展示压力窗口结果。压力窗口从 1M trace 中选择具有实际 eviction 压力的连续区间，训练和评估配置保持不变。该实验是本文最强的真实 workload 正结果：在 \texttt{streamcluster} pressure window 上，QMAP 将加权访问代价从最佳 baseline CLOCK 的 301,767 降至 264,501，降低 12.35\%；迁移次数从 8,937 降至 5,541，降低 38.0\%；命中率从 95.52\% 提升到 97.22\%。这表明，当真实 workload 进入存在明显迁移压力的阶段时，QMAP 能够利用访问历史和候选页特征显著改善迁移决策。

\begin{table*}[t]
  \centering
  \caption{压力窗口实验结果。该实验用于避免标准测试段 eviction 压力不足造成的不可比较问题。}
  \label{tab:pressure-window}
  \scriptsize
  \begin{tabular}{llrrrrr}
    \toprule
    Workload & Policy & Hit rate(\%) & NVM writes & Cost & Migrations & Decision ms \\
    \midrule
    streamcluster & LRU & 95.35 & 585 & 305,562 & 9,276 & 0.000318 \\
    streamcluster & Random & 95.01 & 777 & 314,183 & 9,955 & 0.001190 \\
    streamcluster & LFU & 92.84 & 740 & 361,877 & 14,311 & 0.005936 \\
    streamcluster & CLOCK & 95.52 & \textbf{574} & 301,767 & 8,937 & 0.001553 \\
    streamcluster & QMAP & \textbf{97.22} & 589 & \textbf{264,501} & \textbf{5,541} & 2.360334 \\
    \midrule
    dedup & LRU & \textbf{99.95} & \textbf{99} & \textbf{201,567} & \textbf{87} & 0.000418 \\
    dedup & Random & 99.93 & 105 & 202,076 & 130 & 0.001658 \\
    dedup & LFU & 99.88 & 233 & 203,845 & 221 & 0.006130 \\
    dedup & CLOCK & \textbf{99.95} & 100 & 201,584 & 88 & 0.001643 \\
    dedup & QMAP & \textbf{99.95} & \textbf{99} & \textbf{201,567} & \textbf{87} & 9.091155 \\
    \bottomrule
  \end{tabular}
\end{table*}

\texttt{dedup} pressure window 则是边界案例。QMAP 与 LRU 在命中率、NVM 写入、迁移次数和加权代价上完全持平，而 LFU 和 Random 更差。该结果不能证明 QMAP 在 \texttt{dedup} 上有显著优势，但说明 QMAP 在该弱压力窗口下没有引入额外代价。由于该窗口的迁移次数只有 87 次，策略间可区分空间有限，因此本文不将其作为主要正结果。

\subsection{真实数据消融}

为了分析 QMAP 的收益来源，我们在两个代表性 workload 上做必要消融：\texttt{streamcluster} pressure window 代表强正结果，\texttt{blackscholes} 代表标准切分中的小幅正结果。消融只保留两个最关键变体：去除读写类型特征（no\_rw）和去除写敏感/迁移代价相关损失项（no\_cost）。表~\ref{tab:real-ablation} 给出结果。

\begin{table}[t]
  \centering
  \caption{真实数据消融实验。vs QMAP 表示相对 QMAP 的 cost 变化，正数表示变差。}
  \label{tab:real-ablation}
  \scriptsize
  \begin{tabular}{llrrr}
    \toprule
    Workload & Variant & Cost & vs QMAP & Migrations \\
    \midrule
    streamcluster & QMAP & 264,501 & +0.00\% & 5,541 \\
    streamcluster & no\_rw & 265,905 & +0.53\% & 5,667 \\
    streamcluster & no\_cost & 265,395 & +0.34\% & 5,625 \\
    \midrule
    blackscholes & QMAP & 105,983 & +0.00\% & 525 \\
    blackscholes & no\_rw & 105,961 & -0.02\% & 523 \\
    blackscholes & no\_cost & 105,983 & +0.00\% & 525 \\
    \bottomrule
  \end{tabular}
\end{table}

在 \texttt{streamcluster} pressure window 上，两个消融都会使加权访问代价上升。去除读写类型特征后，cost 上升 0.53\%，迁移次数从 5,541 增至 5,667；去除代价相关损失项后，cost 上升 0.34\%，迁移次数增至 5,625。这说明 QMAP 在主要正结果中确实利用了读写上下文和代价感知训练信号。相比之下，\texttt{blackscholes} 上消融差异接近 0，说明该 workload 的工作集较小，策略收益主要来自少量迁移决策，消融不容易暴露明显机制差异。

\subsection{多随机种子稳定性}

最后，我们验证 QMAP 的结果是否只是单个训练随机种子的偶然现象。我们选择三个 workload：\texttt{streamcluster} pressure window 是最强正结果，\texttt{blackscholes} 是标准切分正结果，\texttt{canneal} 是稳定负例。每个 workload 使用三个训练种子：3136859、42 和 2026。baseline 固定复用确定性或固定随机种子结果，只重新训练 QMAP。

\begin{table*}[t]
  \centering
  \caption{多随机种子稳定性。Delta 表示 QMAP 相对最佳 baseline 的 cost 变化。}
  \label{tab:seed-stability}
  \scriptsize
  \begin{tabular}{lrrrrr}
    \toprule
    Workload & Mean delta & Std delta & Min delta & Max delta & 结论 \\
    \midrule
    streamcluster pressure & -11.59\% & 0.83\% & -12.35\% & -10.43\% & 三个种子均优于最佳 baseline \\
    blackscholes & -1.61\% & 0.51\% & -2.10\% & -0.91\% & 三个种子均优于最佳 baseline \\
    canneal & +18.86\% & 1.63\% & +16.67\% & +20.59\% & 三个种子均差于最佳 baseline \\
    \bottomrule
  \end{tabular}
\end{table*}

表~\ref{tab:seed-stability} 显示，\texttt{streamcluster} pressure window 的平均改善为 11.59\%，标准差为 0.83\%，三个种子均明显优于最佳 baseline。 \texttt{blackscholes} 的改善幅度较小，平均为 1.61\%，但三个种子也全部优于 LFU。 \texttt{canneal} 的失败同样稳定，三个种子均显著差于 LRU。这组结果支持一个平衡结论：QMAP 的主要正结果不是训练 seed 偶然造成的；同时，QMAP 在 \texttt{canneal} 上的失败也是稳定边界，而非单次训练噪声。

\subsection{实验结论}

综合以上结果，QMAP 在真实 trace 上表现出明确但有边界的收益。第一，在具有实际迁移压力的 \texttt{streamcluster} pressure window 上，QMAP 相比最佳 baseline 降低 12.35\% 的加权访问代价，并显著减少迁移次数，这是本文最强的真实 workload 正结果。第二，在 \texttt{blackscholes} 标准 1M split 上，QMAP 稳定小幅优于 LFU，并显著减少 NVM 写入，说明即使在高命中率、小工作集场景中，面向代价的迁移决策仍然可能带来收益。第三，\texttt{dedup} pressure window 上 QMAP 与 LRU 持平，说明在迁移压力较弱或策略空间较小时，QMAP 不一定产生额外优势。第四，\texttt{canneal} 是稳定负例，QMAP 的迁移次数明显增加，暴露出当前候选页排序对特定访问模式仍然敏感。

因此，本文不把 QMAP 表述为所有 workload 上都优于经典策略的通用替换器，而将其定位为面向 DRAM/NVM 混合内存、能在有明显迁移压力和访问相位特征的真实 workload 中降低加权访问代价的学习型页面迁移策略。当前结果也指出后续优化方向：需要进一步增强候选页排序的保守性和在线自适应能力，以避免在 \texttt{canneal} 这类 workload 上出现过度迁移。
